% Auto-generated by scripts/06_emit_structure.py
% Source: scans 018–019
% Section id: frontmatter/preamble
% Phase 3 deliverable. Footnotes (Phase 4), citations (Phase 5),
% and Wade-Giles → Pinyin (Phase 6) are not yet applied here.

\chapter*{Preamble}
\label{fm:preamble}

% source: scan 018, printed 1
The sun's rays glint first on mountains to the west, then, moments later, touch
the thatched roofs of the temples and pit dwellings that follow the curve of the Huan.
The river, still in shadow at the foot of the earthen cliff, winds to the southeast between
clearings of sprouting millet, on its way to merge with the powerful Ho. The year is the
eleventh of Wu Ting's reign, the season spring, the day hsin-wei, eighth of the week.
Filtering through the portal of the ancestral temple, the sunlight wakens the eyes of
the monster mask, bulging with life on the garish bronze tripod. At the center of the temple
stands the king, at the center of the four quarters, the center of the Shang world. Ripening
millet glimpsed through the doorway shows his harvest rituals have found favor. Bronze
cauldrons with their cooked meat offerings invite the presence of his ancestors, their bodies
buried deep and safely across the river, but their spirits, some benevolent, some not, still
reigning over the royal house and the king's person. One is angry, for the king's jaw ached
all the night, is aching now, on the eve of his departure to follow Chih Kuo on campaign
against the 八方.
Five turtle shells lie on the rammed-earth altar. The plastrons have been polished
like jade, but are scarred on their inner side with rows of oval hollows, some already blackened by fire. Into one of the unburned hollows, on the right side of the shell, the diviner
Ch'üeh is thrusting a brand of flaming thorn. As he does so, he cries aloud, ``The sick tooth
is not due to Father Chia!'' Fanned by an assistant to keep the glowing tip intensely hot,
the stick flames against the surface of the shell. Smoke rises. The seconds slowly pass. The
stench of scorched bone mingles with the aroma of millet wine scattered in libation. And
then, with a sharp, clear, puklike sound, the turtle, most silent of creatures, speaks. A
pu-shaped crack has formed in the hollow where the plastron was scorched. Once again
the brand is thrust, now into a matching hollow on the left side of the shell: ``It is due to
Father Chia!'' More time passes . . . another crack forms in response. Moving to the next
plastron, Ch'üeh repeats the charges: ``It is not due to Father Chia!'' Puk! ``It is due to
Father Chia!'' He rams the brand into the hollows and cracks the second turtle shell, then
the third, the fourth, and the fifth.
The diviners consult. The congregation of kinsmen strains to catch their words, for
the curse of a dead father may, in the king's eyes, be the work of a living son. Ch'üeh rubs
wood ash from the fire into the new set of cracks and scrutinizes them once more. But the
shell has given no indication. The charge must be divined again. Two more cracks are
made in each of the five plastrons... and there is again no sign.

% source: scan 019, printed 2
Another brand is plucked from the fire and the new charge cried: ``The sick tooth is
not due to Father Keng!. It is due to Father Keng.'' Father Keng—the king's senior
uncle. This time the indications are clear. His sons, the king's older cousins, turn away
in dismay at the diviner's reading of the cracks. The spirit, their father, has been blamed.
But still the work of spiritual identification continues. ``It is not due to Father Hsin! ….. It
is due to Father Hsin!'' Ch'üeh moves methodically down the row of five plastrons, reciting
the negative and positive charges and cracking each shell twice in this way. No judgment
can be made. Once again, as for Father Chia, ten more cracks are burned. “Auspicious!”
Ch'üeh points to two cracks on the second and fourth shells. Father Hsin is without blame,
his descendants relieved. ...
Now the king speaks. Assistants drag two victims into the temple. There is the barking
and bleating of animals in panic, then silence. Blood stains the earth floor. The king dismembers the victims as Ch'üeh proposes a new charge: ``We sacrifice a dog to Father Keng,
and butcher a sheep.'' The brand flames... puk ... puk... puk... the plastrons crack in
slow and stately sequence. Has the sacrifice mollified the dead uncle? Will the pain in the
sick tooth depart? The king, his hands still sticky with blood, scans the cracks. . . .
In such an atmosphere and in such ways-in a routine that must have consumed tens
of thousands of hours during the Shang historical period-the Shang kings and their
diviners sought to know and fix the future. As the ceremony ended, the diviner handed
the five plastrons to scribes, who began the task of carving into the shell's smooth front
a record of the charges proposed and results observed.\footnote[1]{This imaginative reconstruction, which the specialist may wish to pass over in decent silence, is generally based on evidence that can be documented in the Shang inscriptions. The plastron set is Ping-pien 12-21, discussed in detail in sec. 3.7. The actual year of Wu Ting's reign is hypothetical; for the season, see ch. 3, nn. 85, 98. For the description of 小屯, I rely on Karlbeck (1957), pp. 174-192; on Tung (1965), pp. 14-15; on the reconstruction of the cult center given in Wheatley (1971), pp. 39-47; and on my own observations during a visit to the site on 1 June 1975. For evidence that some Shang divination took place at daybreak, see Tung (1945), pt. 2, ch. 5, p. 23a. The numerous divination charges of the form ``Today it will not rain” (S159.3) or ``Today the king hunts'' (S293.4) were presumably divined early on the day in question. The Shang, undoubtedly, were early to bed and early to rise (cf. Needham [1959], p. 419, note b). I have arbitrarily selected building Yi Z 8 (whose remains are described in detail by 石璋如 [1959], pp. 89-99) as the scene of the divinations. For evidence that divination took place in the ancestral temple, see, e.g., Chin-chang 120 (D); Ch'ien-pien 8.15.1; Nan-pei, “Ming'' 729 (D) (all S270.2); and similar inscriptions at S270.3. For the possibility that this set was divined on campaign and not in the ritual center at 小屯, see ch. 3, n. 98. On the proverbial silence of the turtle, see Pope (1955), p. 66; on the ``voice of the turtle,'' see ch. 1, n. 97. For the social and political implications of divination, see the works cited in ch. 5, n. 8. Exactly how the pyromantic cracks were read and by whom is not entirely clear (see secs. 2.6 and 2.7; ch. 2, n. 58). Whether it was the king or an assistant who actually dismembered the victims is not specified, but the king is recorded as sacrificer in other cases (e.g., Chia-pien 182; Yi-pien 9054, 9103; and other inscriptions at \$507.1); in any event, sacrifice divinations were intimately linked to the shedding of blood.}¹
The Shang kings read the mantic cracks to divine the wishes of their ancestors. We
read the mantic inscriptions to divine the wishes of the Shang kings. May the oracle bones,
once used to read the future, now be used to read the past!
